\section{Two-symbol alphabets}\label{sec:real}

The alphabets of Peres and of Conway and Kochen are both of the form $\{0,\pm1,\pm x\}$.  For any
$x\notin\{0,\pm1\}$ this alphabet has $5^3-1=124$ words, the same list for every~$x$
(Fig.~\ref{fig:generic}).  They give 49 rays for every $x$ other than $\pm i$, where more
words coincide and 31 remain (Theorem~\ref{thm:twosymbol}).  What depends on $x$ is which
inner products vanish.

\subsection{Which inner products can vanish}

A term of the inner product of two words is $a^{*}b$ with $a,b$ letters, so it is zero or one of
$\pm1,\ \pm x,\ \pm x^{*},\ \pm|x|^2$.  An inner product is therefore
\[
c_0+c_1x+c_2x^{*}+c_3|x|^2,\qquad c_i\in\Z,\quad |c_0|+|c_1|+|c_2|+|c_3|\le 3 .
\]
Some of these vanish for every $x$ (all $c_i=0$: disjoint supports, or terms cancelling in
pairs).  The others vanish only for special~$x$, and since there are finitely many coefficient
vectors they can be listed completely.  There are 128 nonzero vectors.  The pool is unchanged by
$x\mapsto-x$ and by $x\mapsto1/x$, and complex conjugation of all coordinates carries it to the
pool of $x^{*}$, which has the same orthogonality graph.  Together with an overall sign these
maps act on the coefficient vectors, inversion as
$(c_0,c_1,c_2,c_3)\mapsto(c_3,c_2,c_1,c_0)$.  The 128 vectors fall into 21 orbits, of which 13
have no solution other than $x\in\{0,\pm1\}$.  Table~\ref{tab:zerosums} lists the remaining
eight.

\begin{table}[t]
\centering
\caption{Every relation of at most three terms that can make an inner product of two words over
$\{0,\pm1,\pm x\}$ vanish for some but not all $x$, up to the symmetries $x\mapsto-x$,
$x\mapsto x^{*}$, $x\mapsto1/x$ and overall sign (computed by \texttt{ks\_zero\_sums\_v10.py}).
The row $x-x^{*}=0$ is the reason real alphabets have orthogonalities that generic complex ones
lack.
The rows $2+x=0$ and $1+2x=0$ describe the same pool, since $S$ is unchanged by $x\mapsto 1/x$.
A row says where extra orthogonalities appear; it says nothing by itself about colorability.}
\label{tab:zerosums}
\input{figures/zero_sums_table.tex}
\end{table}

Three remarks.  First, the solution sets are not isolated points in general: two of them are
circles and three are lines in the complex plane.  Second, the two-term relations were missing
from the corresponding table of earlier versions; the relation $x+x^{*}=0$ was pointed out to us
by Frembs~\cite{Frembs2026}.  It holds on the imaginary axis, as $x-x^{*}=0$ holds on the real
axis, and the two classical modulus-2 alphabets differ in exactly this.  Of the 120 orthogonal
pairs of the $\sqrt{-2}$ pool (in the word representatives used here), 60 hold for every $x$, 54
use $x+x^{*}=0$ and 6 use $|x|^2=2$; at $x=\sqrt2$ the 54 are replaced by pairs using
$x-x^{*}=0$.  The difference is visible in the graphs (Computed fact~\ref{fact:pools}).  Third, the statement that an
uncolorable pool must have some vanishing inner products of this kind is true by definition and
has no content; versions 1--9 presented it as a finding.  The content is in which loci give
uncolorable pools, and that has to be computed.

\subsection{Real alphabets}

\begin{table}[t]
\centering
\caption{Raw pools of two-symbol alphabets $\{0,\pm1,\pm x\}$, exact arithmetic.  ``Min'' is the
certified size of the smallest uncolorable subset.  Every real $x=\sqrt d$ with $d$ a non-square
integer, $3\le d\le30$, gives the same counts as $\sqrt3$, as do $x=3$, $x=4$ and
$x=(1+\sqrt d)/2$ for $d=13,17,21,29$.}
\label{tab:alphabets}
\begin{tabular}{lrrrcc}
\toprule
Alphabet & Rays & Pairs & Triads & Colorable & Min \\
\midrule
$\{0,\pm1\}$ & 13 & 24 & 4 & yes & --- \\
$\{0,\pm1,\pm2\}$ (also $\pm\tfrac12$) & 49 & 138 & 26 & \textbf{no} & \textbf{31} \\
$\{0,\pm1,\pm\sqrt2\}$ & 49 & 120 & 16 & \textbf{no} & \textbf{33} \\
$\{0,\pm1,\pm\sqrt3\}$ & 49 & 114 & 10 & yes & --- \\
$\{0,\pm1,\pm\varphi\}$, $\varphi=(1+\sqrt5)/2$ & 49 & 138 & 10 & yes & --- \\
\midrule
$\{0,\pm1,\pm\sqrt{-1}\}$ & 31 & 69 & 7 & yes & --- \\
$\{0,\pm1,\pm\sqrt{-2}\}$ & 49 & 120 & 16 & \textbf{no} & \textbf{33} \\
$\{0,\pm1,\pm\sqrt{-d}\}$, $d=3,5,6,\dots,15$ & 49 & 114 & 10 & yes & --- \\
\bottomrule
\end{tabular}
\end{table}

Table~\ref{tab:alphabets} gives the raw pools.  The 13-ray pool of $\{0,\pm1\}$ is colorable, as
it must be: 13 is below the lower bound of~24.  Among real two-symbol alphabets the
only uncolorable raw pools are those of $x=2$ and $x=\sqrt2$ (and their negatives and
reciprocals), the two alphabets of Peres's cubes; Theorem~\ref{thm:twosymbol} below proves this
for every real~$x$ and settles the complex case as well.
The row for $\sqrt{-1}$ corrects version~9, which printed 25 rays; every nonzero letter of that
alphabet is a unit, the group $\{\pm1,\pm i\}$ acts freely on the 124 words, and $124/4=31$.

\begin{fact}[Integer pool]\label{fact:integer}
The 49-ray pool of $\{0,\pm1,\pm2\}$ has minimum uncolorable subsets of 31 rays, certified.  There
are exactly six such subsets.  They form a single orbit under the 48 signed permutations of the
coordinates, so all six are copies of $\CK$.  Together they use 37 of the 49 rays; the 12 unused
rays are those of squared length~9, such as $(1,2,2)$; 13 rays, those of $\{0,\pm1\}$, lie in all
six.  The completed pool has 109 rays, 258 pairs and 86 triads, and its minimum is also 31,
certified.
\end{fact}

Version~9 stated that the pool contains 162 distinct minimal sets falling into six classes.  The
exhaustive enumeration gives six sets in one class.  The pool also contains uncolorable subsets
that are critical (no ray can be removed) but not of minimum size, among them the 33-ray set
of Sch\"utte, also known as the Conway--Kochen 33-ray set and communicated by them to
Peres~\cite{Peres1991} (the identification is made in~\cite{TrandafirCabello2025}), which has 76
orthogonal pairs and
20 bases and is contained in the 37 rays above.

None of three enlargements we tried lowers the minimum: the raw pools of $\{0,\pm1,\pm2,\pm3\}$ (145 rays),
$\{0,\pm1,\pm\sqrt2,\pm2\}$ (109 rays) and $\{0,\pm1,\pm\varphi,\pm2\}$ (145 rays) all have
certified minimum 31, with 17 bases in the minimum set found, while
$\{0,\pm1,\pm\varphi,\pm\varphi^2\}$ (109 rays, 24 triads) is colorable.  Each of the first three
contains the integer pool, so 31 is an upper bound by Lemma~\ref{lem:restrict}; the content of
the computation is that nothing smaller appears.

\begin{fact}[Peres pool]\label{fact:peres}
The 49-ray pool of $\{0,\pm1,\pm\sqrt2\}$ has a unique minimum uncolorable subset, the 33 rays of
Peres (72 pairs, 16 bases).  The other 16 rays lie in no triad.  The completed pool has 145 rays,
336 pairs and 112 triads, and its minimum is 33, certified.  The same statements hold for
$\{0,\pm1,\pm\sqrt{-2}\}$.
\end{fact}

\begin{fact}[The two modulus-2 pools are not the same graph]\label{fact:pools}
The minimum 33-ray sets of the $\sqrt2$ and $\sqrt{-2}$ pools have isomorphic orthogonality
graphs.  The two 49-ray pools do not, although they agree in the number of rays, pairs, triads,
in degree sequence and in spectrum; neither do the two 145-ray completed pools.  (Graph
isomorphism was tested with VF2~\cite{Cordella2004VF2}; for the raw pools we also checked directly
that none of the 48 isomorphisms between the two 33-ray cores carries the neighbourhoods of the
other 16 rays of one pool onto those of the other.)
\end{fact}

Versions 1--9 called these two alphabets ``isomorphic islands''.  That is correct for the
minimum sets and false for the pools.  The difference is the one Frembs identified: 54 of the 120
orthogonalities of the $\sqrt{-2}$ pool use $x+x^{*}=0$, a relation that fails at $x=\sqrt2$, and
he builds an arrangement on that relation~\cite{Frembs2026}.

\subsection{All two-symbol alphabets}\label{sec:twosymbol}

The completeness of Table~\ref{tab:zerosums} turns the question ``for which $x$ is the pool of
$\{0,\pm1,\pm x\}$ uncolorable?'' into a finite computation, and the answer is a theorem about
every complex~$x$.

\begin{theorem}[Two-symbol alphabets]\label{thm:twosymbol}
Let $x$ be a nonzero complex number and $S(x)$ the pool of all words over $\{0,\pm1,\pm x\}$.
\begin{enumerate}[label=(\roman*)]
\item $S(x)$ is uncolorable if and only if $x$ or $1/x$ is one of $\pm2$, $\pm\sqrt2$,
$\pm\sqrt{-2}$, or $x$ is a primitive third or sixth root of unity.
\item The minimum uncolorable subset of $S(x)$ has 31 rays for $x\in\{\pm2,\pm\tfrac12\}$, where
the minimum sets are the six copies of $\CK$ of Computed fact~\ref{fact:integer}, and 33 rays in
every other uncolorable case.
\item Consequently, if all rays of a KS set have coordinates in $\{0,\pm a,\pm b\}$ for some
complex numbers $a$ and $b$, the set has at least 31 rays.
\end{enumerate}
\end{theorem}

\begin{proof}
Write $x=a+ib$.  The orthogonality graph of $S(x)$ depends on $x$ only through two finite pieces
of data: which words are proportional, and which inner products vanish.

\emph{Proportionality.}  Two words are proportional when their three $2\times2$ minors vanish.
A minor is a difference of two products of two letters, and such a product is $0$, $\pm1$, $\pm x$ or $\pm x^2$.  So a minor is identically zero, or a
monomial or twice a monomial (never zero), or one of $\pm(1\pm x)$, $\pm(1\pm x^2)$,
$\pm x(1\pm x)$.  Away from $x\in\{\pm1,\pm i\}$
the same words are proportional for every~$x$, and the pool has 49 rays.

\emph{Orthogonality.}  An inner product is $c_0+c_1x+c_2x^{*}+c_3|x|^2$ for one of the 128
coefficient vectors counted above, or is identically zero.  It vanishes when
$(c_1-c_2)\,b=0$ and $c_0+(c_1+c_2)\,a+c_3(a^2+b^2)=0$.  For a given vector the solution set is
therefore empty, a finite set of real points, the real axis, a vertical line, or a circle
centred on the real axis.  Going through the 128 vectors gives nine curves,
\[
\mathrm{Im}\,x=0,\quad \mathrm{Re}\,x=0,\ \pm\tfrac12,\quad |x|^2=\tfrac12,\ 1,\ 2,\quad
|x\pm1|=1,
\]
and finitely many real points.  Call a point \emph{special} if it is one of those real points, a
point of intersection of two of the curves, or one of $\pm1,\pm i$.  There are 27 nonzero special points
up to complex conjugation.  At a point that is not special and lies on none of the curves, no
vector vanishes.  At a point of a curve $C$ that is not special, the vectors that vanish are
exactly those whose solution set is~$C$, since any other vector vanishes on a different curve or
on special points only.  Hence the graph of $S(x)$ is constant on the complement of the curves,
constant on each curve away from the special points, and otherwise has to be computed at each
special point: $1+9+27=37$ cases.  (Conjugating all coordinates carries $S(x)$ to $S(x^{*})$ and
preserves orthogonality, so $x$ and $x^{*}$ need not be distinguished.)

\emph{Computation.}  The script \texttt{ks\_two\_symbol\_theorem.py} lists the curves and special
points, takes an exact sample point on each curve, builds the 37 graphs in exact arithmetic, decides colorability of each, and runs the certification procedure of
Section~\ref{sec:methods} on every uncolorable one; all terminate.
Table~\ref{tab:twosymbol} is the result.  Special points are grouped by the symmetries $x\mapsto-x$
and $x\mapsto1/x$, which do not change the pool; curves are listed on one row when their counts
agree, which does not mean their graphs are isomorphic (the generic graphs of the real axis, the
imaginary axis and the unit circle are pairwise non-isomorphic).  Parts (i) and (ii) are read off the table.  For
(iii), if $a$ or $b$ is zero or $a=\pm b$ the rays lie in the 13-ray pool of $\{0,\pm1\}$, which
is colorable; otherwise dividing all coordinates by $a$ puts the set inside $S(b/a)$.
\end{proof}

\begin{table}[t]
\centering
\caption{Every pool of a two-symbol alphabet $\{0,\pm1,\pm x\}$, $x\in\mathbb{C}\setminus\{0\}$
(Theorem~\ref{thm:twosymbol}).  A curve stands for its points other than the special points
listed below it.  ``Min'' is the certified size of the smallest uncolorable subset, and the last
column names the graph of the minimum set found (isomorphism by VF2).}
\label{tab:twosymbol}
\small
\begin{tabular}{@{}lrrrcl@{}}
\toprule
$x$ & Rays & Pairs & Triads & Min & Minimum set \\
\midrule
on no curve & 49 & 60 & 4 & --- & colorable \\
$\mathrm{Im}\,x=0$;\ \ $\mathrm{Re}\,x=0$;\ \ $|x|=1$ & 49 & 114 & 10 & --- & colorable \\
$\mathrm{Re}\,x=\pm\tfrac12$;\ \ $|x\pm1|=1$ & 49 & 72 & 8 & --- & colorable \\
$|x|^2=2$;\ \ $|x|^2=\tfrac12$ & 49 & 66 & 10 & --- & colorable \\
\midrule
$\pm1$ & 13 & 24 & 4 & --- & colorable \\
$\pm i$ & 31 & 69 & 7 & --- & colorable \\
$\pm\varphi,\ \pm1/\varphi$ & 49 & 138 & 10 & --- & colorable \\
$\pm1\pm i,\ (\pm1\pm i)/2$ & 49 & 78 & 14 & --- & colorable \\
$(\pm1\pm\sqrt{-7})/2,\ (\pm1\pm\sqrt{-7})/4$ & 49 & 78 & 14 & --- & colorable \\
$\pm2,\ \pm\tfrac12$ & 49 & 138 & 26 & \textbf{31} & $\CK$ \\
$\pm\sqrt2,\ \pm1/\sqrt2$ & 49 & 120 & 16 & \textbf{33} & Peres \\
$\pm\sqrt{-2},\ \pm1/\sqrt{-2}$ & 49 & 120 & 16 & \textbf{33} & Peres \\
$(\pm1\pm\sqrt{-3})/2$ & 49 & 138 & 18 & \textbf{33} & Eisenstein 33 \\
\bottomrule
\end{tabular}
\end{table}

Four remarks on what the theorem does and does not say.

First, it is a statement about a class of coordinate systems, not about KS sets in general.
The minimum number of rays of a KS set in dimension three is known only to lie between 24
and~31, and a set smaller than 31, if one exists, would need at least three nonzero letters up to
sign.  What the theorem
shows is that the construction read off Peres's cubes, one unit and one further number along
each axis, cannot beat Conway and Kochen whatever the number is.

Second, it concerns the raw pool.  Completion adds rays whose coordinates are not letters, and
the completed pool of $\{0,\pm1,\pm\varphi\}$ is uncolorable although its raw pool is not
(Section~\ref{sec:catalogue}).

Third, the uncolorable cases are isolated points; no curve of Table~\ref{tab:zerosums} is
uncolorable for a two-symbol alphabet.  In particular the two-symbol pool is colorable at a
generic point of the circle $|x|^2=2$, and even at $x=1+i$ and $x=(1+\sqrt{-7})/2$, where it
meets another curve.  The uncolorability of that whole circle (Theorem~\ref{thm:circle}) needs the
conjugate letter $x^{*}$.

Fourth, the last row did not appear in earlier versions.  At a primitive sixth root of unity three curves
meet ($|x|=1$, $\mathrm{Re}\,x=\tfrac12$, $|x-1|=1$), and the 49 rays over $\{0,\pm1,\pm\omega\}$
are already uncolorable, with a minimum set isomorphic to the Eisenstein 33-ray set of
Section~\ref{sec:roots}.  The letter $\omega^2$ of the Eisenstein alphabet is not needed.  We do
not claim this set is new: it has the graph of Cabello's 33-ray set~\cite{Cabello2025simplest}.

Here ``coordinates in an alphabet'' means, as throughout, that every ray has a representative
whose three entries are letters, in one common orthonormal basis.  We have not compared
Theorem~\ref{thm:twosymbol} with the discussion of alphabets in the preprints of Flores
Gordillo~\cite{FloresGordillo2026} and make no claim of priority for the classification.

\subsection{Real three-symbol alphabets}\label{sec:threesymbol}

The same method works with one more real letter.  Signs are in the alphabet and dividing every
word by a letter does not change the pool, so a real alphabet $\{0,\pm a,\pm b,\pm c\}$ whose
letters have three distinct nonzero absolute values may be written $\{0,\pm1,\pm x,\pm y\}$ with
$1<x<y$; if two absolute values agree or a letter is zero, the alphabet has at most two symbols and
Theorem~\ref{thm:twosymbol} applies.

\begin{theorem}[Real three-symbol alphabets]\label{thm:threesymbol}
Let $1<x<y$ be real and $S(x,y)$ the pool of all words over $\{0,\pm1,\pm x,\pm y\}$.
\begin{enumerate}[label=(\roman*)]
\item $S(x,y)$ is uncolorable if and only if two of the letters $1,x,y$ have ratio $2$ or
$\sqrt2$, or $(x,y)$ is one of the eight points of Table~\ref{tab:threesymbol}.
\item Its minimum uncolorable subset has 31 rays when two letters have ratio~2, 33 rays when two
letters have ratio $\sqrt2$ and none has ratio~2, and between 42 and 57 rays at the eight points
(Table~\ref{tab:threesymbol}).
\item Consequently, if all rays of a KS set have coordinates in $\{0,\pm a,\pm b,\pm c\}$ for some
real numbers $a,b,c$, the set has at least 31 rays.
\end{enumerate}
\end{theorem}

\begin{proof}
An inner product of two words, and a $2\times2$ minor of two words, is
$p_c=c_0+c_1x+c_2y+c_3x^2+c_4xy+c_5y^2$ for an integer vector~$c$ with
$\sum_k|c_k|\le3$; 376 nonzero vectors occur, counting $c$ and $-c$ separately.  Factor each $p_c$ over~$\Q$;
125 distinct irreducible factors appear.  Since $p_c(P)=0$ exactly when one of its factors vanishes
at~$P$, which words are proportional, and hence which words represent the same ray, and which
are orthogonal depend only on the set $F(P)$ of irreducible factors that vanish at~$P$; so
does the orthogonality graph.  The sets $F(P)$ that occur for $P$ in the domain $D=\{1<x<y\}$ are:
the empty set; $\{f\}$ for each factor $f$ whose zero set meets $D$ in a curve (31 of them); and
the set $F(P)$ at each point of $D$ that lies on two of these curves or is a singular point of
one.  There are 130 such special sets.  The script \texttt{ks\_three\_symbol\_real.py} makes
every one of these decisions exactly, holding real algebraic numbers as rational isolating
intervals and deciding signs by exact interval arithmetic.  It finds the curves by cutting the
line $x>1$ at the real roots of the leading coefficient and discriminant in~$y$ and of
$f(x,x)$: on each cell the real branches over~$x$ neither meet nor cross $y=x$, so one rational
point per cell decides whether a branch lies in~$D$.  It finds the special points by splitting
each finite intersection, and the singular locus of each curve, into prime components
$\{(h_x(t),h_y(t)) : q(t)=0\}$ with $q$ irreducible, using a separating coordinate
$t=x+sy$ after making the ideal radical.  A factor $h$ vanishes on a component exactly
when $q$ divides $h(h_x,h_y)$.  It then builds the $1+31+130=162$ graphs and decides
colorability.  As a cross-check it evaluates every special point to 60 digits.  It also
rebuilds the graph (the partition of the words into rays and the set of orthogonal pairs) at a
point of each of the 162 strata, this time from the coordinates and without the factorisation,
and finds it identical in every case.  An earlier version of the scan that made the same
decisions in floating point gave the same 162 cases.

Six curves are uncolorable: $x=2$, $y=2$, $y=2x$, $x^2=2$, $y^2=2$, $y^2=2x^2$.  The generic
point, the other 25 curves and 65 of the special sets are colorable; the other 65 special sets
are uncolorable, and all but eight of them lie on one of the six curves.  Every uncolorable case
was certified (Section~\ref{sec:methods}), each run also proved separately that no
uncolorable subset has fewer than 31 rays, and a minimum set for each case is archived and
checked uncolorable.  Points of the six curves have the minimum of the
curve, except that a point on a ratio-2 curve and a ratio-$\sqrt2$ curve has minimum~31.  This
gives (i) and~(ii).  For (iii), if two of $|a|,|b|,|c|$ agree or one is zero, the set lies in a
two-symbol pool and Theorem~\ref{thm:twosymbol}(iii) applies; otherwise, dividing by the smallest
absolute value and renaming puts it inside some $S(x,y)$.
\end{proof}

\begin{table}[t]
\centering
\caption{The eight points of Theorem~\ref{thm:threesymbol} that lie on none of the six
uncolorable curves, with the relations that hold there (each point is the unique point of $D$
satisfying them) and the certified minimum.  Rays, pairs and triads are those of the raw pool.}
\label{tab:threesymbol}
\small
\begin{tabular}{@{}llrrrr@{}}
\toprule
Relations & $(x,y)\approx$ & Rays & Pairs & Triads & Min \\
\midrule
$x^2=y,\ y^2=y+1$\ \ ($x=\sqrt\varphi,\ y=\varphi$) & $(1.272,1.618)$ & 109 & 312 & 28 & 42 \\
$x^3=x^2+1,\ y=x^2$\ \ ($x=\psi,\ y=\psi^2$) & $(1.466,2.148)$ & 109 & 348 & 40 & 47 \\
$y^2=x^2+1,\ xy=x+y$ & $(1.883,2.132)$ & 145 & 474 & 42 & 49 \\
$y^2=x^2+1,\ x^2=2y$\ \ ($y=1+\sqrt2$) & $(2.197,2.414)$ & 145 & 402 & 46 & 49 \\
$xy=x^2+1,\ y=x^2-x,\ y^2=x^2+x$ & $(2.206,2.659)$ & 145 & 462 & 46 & 51 \\
$xy=x^2+1,\ x^2=2y$ & $(2.359,2.783)$ & 145 & 414 & 58 & 55 \\
$y=x+1,\ x^2=2y$\ \ ($x=1+\sqrt3$) & $(2.732,3.732)$ & 145 & 486 & 46 & 56 \\
$xy=2,\ y^2=x^2+1$ & $(1.250,1.600)$ & 145 & 390 & 46 & 57 \\
\bottomrule
\end{tabular}
\end{table}

Three remarks.  First, the curves contribute nothing new: on a ratio-2 curve the pool contains
the words over $\{0,\pm1,\pm2\}$ (after rescaling) and hence $\CK$, and on a ratio-$\sqrt2$ curve
it contains Peres's set; the theorem says that no third letter, placed anywhere on these curves,
gives a smaller set.  Second, the eight points are exceptional uncolorable alphabets (we have
not searched the literature for them), and all their minima exceed~33.  One of them is the supergolden alphabet of Section~\ref{sec:catalogue}, whose
minimum 47 was certified there independently.  Third, the theorem is about real letters; the
complex three-symbol case has the same finite structure in principle but a parameter space of four
real dimensions, and we have not attempted it.
